\documentclass[11pt]{article}
\usepackage[margin=2.4cm]{geometry}
\usepackage{graphicx}
\usepackage{booktabs}
\usepackage{hyperref}
\usepackage[spanish]{babel}
\hypersetup{colorlinks=true, linkcolor=blue, urlcolor=blue}

\title{Track B: resiliencia, buffers y arquitecturas\\
\large Resumen corto y directo para Garrido}
\author{Equipo SCRES-IA}
\date{3 de julio de 2026}

\begin{document}
\maketitle

\section*{De qué se trata}

Este documento responde cuatro cosas que quedaron abiertas en la reunión del 2 de julio: cómo se
reabastecen los buffers, si usar variables continuas contradice el DES, cuál es exactamente el
contrato Track B 8D, y qué papel cumple el \emph{forecast}. Para no distraernos, aquí nos
concentramos casi por completo en resiliencia: usamos la fórmula Excel de Garrido (ReT) como
métrica principal. Las otras métricas sirven para auditar, pero no son el centro de la conversación.

\section*{Buffers: ¿de dónde sale el inventario?}

Tenías razón en dudar. El mecanismo viejo literalmente creaba inventario de la nada (una función
interna que añadía unidades directas al contenedor de inventario), sin descontarlo de ningún
origen ni chequear capacidad. Lo arreglamos: ahora el agente solo mueve los parámetros normales de
reposición (cantidad de orden y punto de reorden) que el simulador ya usa para producir y
transportar de verdad --- todo pasa por la cadena real, con capacidad y tiempo de transporte.

\begin{figure}[htbp]
\centering
\includegraphics[width=0.95\textwidth]{figures/fig3_buffer_mechanism.pdf}
\caption{Mecanismo viejo (el problema) vs. corregido (conserva el flujo físico real).}
\end{figure}

\section*{¿Continuo contradice el DES?}

No. El DES describe cómo avanza el reloj de la simulación (por eventos, no por pasos fijos); no
dice nada sobre si las variables de decisión deben ser discretas. La propia tesis ya usa
fracciones continuas de buffer, y tú mismo recomendaste continuas sobre discretas en la reunión.
No hay conflicto.

\section*{Qué decide el agente}

En Track B no dejamos que la red cambie toda la cadena. Le damos pocas palancas, escogidas porque
están cerca del cuello de botella donde sí puede mejorar la resiliencia. El contrato usado se llama
\textbf{\texttt{track\_b\_v1}}. Es un contrato de \textbf{ocho decisiones}; por eso hablamos de
Track B 8D. Esto no debe confundirse con \textbf{v7}, que es la versión de observación, no una
acción de siete dimensiones.

La acción 8D tiene ocho señales continuas entre $-1$ y $1$:

\begin{itemize}
  \item \textbf{CDC / Op3}: cantidad de reposición y punto de reorden del centro de distribución.
        Esto permite mover materia prima hacia la cadena, pero siempre respetando disponibilidad y
        capacidad. Lo incluimos porque Op3 ya aparece en la tesis como nodo de buffer/inventario;
        además, cuando decimos ``del CDC hacia abajo'', la lectura natural es incluir el CDC como
        nodo frontera y excluir sólo Op1/Op2 (contratación/proveedores).
  \item \textbf{Op9}: cantidad de despacho y punto de reorden hacia el batallón. Es la primera
        palanca claramente aguas abajo del CDC.
  \item \textbf{Op5}: señal heredada del contrato de tesis para el flujo hacia ensamblaje. No es la
        palanca que usamos para defender el resultado principal.
  \item \textbf{Turnos}: una señal continua que se traduce a 1, 2 o 3 turnos de ensamblaje.
  \item \textbf{Op10 y Op12}: multiplicadores de despacho aguas abajo. Estas son las palancas más
        importantes de Track B: regulan cuánto empuja la cadena hacia CSSU y teatro cuando hay
        presión de servicio.
\end{itemize}

Por eso, la versión \textbf{8D completo} es válida y alineada con la tesis: incluye variables ya
familiares (Op3, Op9 y turnos) más la extensión nueva aguas abajo (Op10/Op12). Como Op3 es el CDC,
la frase ``del CDC hacia abajo'' la estamos interpretando de forma inclusiva: incluye Op3 y excluye
Op1/Op2. Esa lectura también es consistente con la tesis, porque Op3 ya estaba en el contrato de
inventario/buffer.

\section*{Qué observa el agente}

La política no ve el futuro ni las órdenes individuales completas. Observa un resumen semanal del
estado de la simulación. En las corridas principales usamos el contrato \textbf{v7}, con 52 señales
normalizadas. No hace falta memorizar las 52; se agrupan así:

\begin{itemize}
  \item \textbf{Inventario disponible}: materia prima en CDC y línea de ensamblaje, raciones en
        ensamblaje, batallón, CSSU y teatro.
  \item \textbf{Servicio y backlog}: fill rate, backorder rate, demanda y backorders de la semana
        anterior, backlog acumulado y edad del backlog.
  \item \textbf{Estado operativo}: si ensamblaje u operaciones logísticas están caídas, turnos
        activos, lote pendiente y demanda contingente.
  \item \textbf{Calendario y régimen}: fase de ciclos de Op1/Op2, día/semana de trabajo, y estado
        del régimen de riesgo (normal, tensionado, pre-disrupción, disrupción, recuperación).
  \item \textbf{Forecast y presión de riesgo}: forecast a 48h y 168h, deuda de mantenimiento,
        probabilidad de defectos R14 y señales recientes de riesgos aguas abajo.
  \item \textbf{Cuello de botella Track B}: estado de Op10/Op12, presión de cola en esos nodos,
        fill rate y backlog móviles de 4 semanas.
\end{itemize}

¿Son demasiadas? Para una persona, sí sería una lista larga; para PPO+MLP o Real-KAN, 52 variables
es un tamaño moderado. Además, no son variables arbitrarias: cada bloque responde a una pregunta
operativa concreta: cuánto inventario hay, cómo va el servicio, qué tan congestionada está la cola,
si hay riesgo activo o cercano, y dónde está la presión aguas abajo. Probamos también versiones con
historia y memoria; hasta ahora no ganaron a esta observación compacta v7.

\section*{Forecast: qué es y por qué hay que tratarlo con cuidado}

El forecast son dos señales normalizadas que resumen presión de riesgo a \textbf{48 horas} y
\textbf{168 horas}. En el benchmark adaptativo se actualizan cada 48 horas. Técnicamente no son una
bola de cristal: salen del régimen de riesgo del simulador, de sus probabilidades de transición, de
la intensidad esperada de disrupción y de un ruido pequeño. En una interpretación operativa, esto
puede verse como una alerta temprana: inteligencia logística, mantenimiento, monitoreo de rutas o
información contextual que anticipa que el sistema entra en una fase más peligrosa.

La ventaja es obvia: si queremos estudiar prevención, el forecast le da al agente una señal para
prepararse antes del golpe. Pero también es polémico: como nace dentro del simulador, un revisor
podría decir que es una señal privilegiada. Si el agente gana sólo porque lee un indicador que un
operador real no tendría, el resultado sería menos defendible. Por eso conviene separarlo en tres
preguntas:

\begin{itemize}
  \item \textbf{¿Lo usamos para entrenar?} Sí, en el contrato v7 completo.
  \item \textbf{¿El modelo depende de él?} Ya hicimos la prueba limpia: reentrenamos Track B sin los
        dos forecasts explícitos. El resultado prácticamente no cambia en ReT Excel: el PPO+MLP
        sin forecast queda en 0.005895, frente a aproximadamente 0.005898 del canónico con forecast.
        Además, sigue ganando con claridad a la mejor estática/heurística comparable.
  \item \textbf{¿Debemos quitarlo?} Para el resultado principal, sí podemos defender una versión sin
        forecast. Es más conservadora y evita la crítica de información privilegiada. Para estudiar
        prevención, en cambio, el forecast sigue siendo útil como sonda: si permite actuar antes del
        riesgo y eso mejora la ReT, entonces se justifica como herramienta de alerta temprana.
\end{itemize}

Mi recomendación para discutir con Garrido: \textbf{no vender forecast como prevención demostrada
todavía}. La buena noticia es que ya no necesitamos forecast para sostener el claim principal:
Track B aprende igual de forma adaptativa y robusta. Dejemos forecast como herramienta para una
pregunta más ambiciosa: si el agente puede volverse preventivo.

\section*{Resiliencia: los resultados}

Aquí muestro sólo el contrato \textbf{8D completo}. Quitamos la variante de CDC congelado del
cuerpo principal porque distrae: era una auditoría útil, pero la propuesta que queremos discutir
con Garrido es el modelo completo, CDC incluido. Probamos dos arquitecturas: PPO+MLP y Real-KAN,
la implementación oficial de \texttt{pykan}.

Todas las comparaciones usan los mismos escenarios de riesgo para cada política (misma semilla,
mismos eventos), así que la diferencia es limpia, no suerte. A eso a veces se le llama CRN
(\emph{common random numbers}), pero en simple significa: \textbf{todos compiten contra el mismo
partido}. El comparador es nuestra propia grilla de políticas estáticas constantes (147
combinaciones de turno y despacho).

La tesis sigue siendo la base: usamos su DES, sus operaciones, sus riesgos y su fórmula Excel de
resiliencia. Lo que no conviene hacer es decir ``mejoramos X\% contra la resiliencia de la tesis'',
porque la tesis usa otro régimen experimental. Para mejora, la comparación justa es contra la mejor
política estática dentro de nuestro mismo experimento.

\begin{figure}[htbp]
\centering
\includegraphics[width=0.85\textwidth]{figures/fig2_results.pdf}
\caption{Mejora de resiliencia (ReT Excel) del contrato 8D completo frente a la mejor política
estática.}
\end{figure}

\textbf{Las dos arquitecturas ganan en el contrato 8D completo.} Estos porcentajes son
la mejora contra nuestra propia grilla estática densa, no contra un número absoluto de la tesis.
Esa es la comparación limpia: mismo régimen, mismos eventos, misma métrica ReT Excel.

\begin{center}
\begin{tabular}{lcc}
\toprule
Alternativa & Superficie de decisión & Mejora vs. mejor estática \\
\midrule
PPO+MLP & 8D completo & +8.0\% \\
Real-KAN & 8D completo & +9.4\% \\
\bottomrule
\end{tabular}
\end{center}

\clearpage

\begin{itemize}
  \item \textbf{Op3/CDC sí puede estar}: es el nodo frontera y ya aparece en la tesis. El contrato
        ampliado no toca Op1/Op2; empieza en el CDC y baja hacia Batallón, CSSU y Teatro.
  \item \textbf{Si miramos sólo resiliencia, Real-KAN es el mejor}: gana un poco más que PPO+MLP
        en el contrato 8D completo. KAN sí funciona; la decisión es si priorizamos máxima
        resiliencia o una arquitectura más simple.
\end{itemize}

\section*{PPO+MLP vs.\ Real-KAN: ¿cuál preferimos?}

Nuestra recomendación para el paper sigue siendo \textbf{PPO+MLP} como columna vertebral: gana de
forma clara, es fácil de defender y no depende de una arquitectura nueva para que el resultado
exista. Pero si el criterio principal de Garrido es \textbf{resiliencia máxima}, entonces Real-KAN
merece una oportunidad real: obtiene la ReT Excel más alta, conecta con la novedad
Kolmogorov-Arnold de su paper de 2024 y tiene una ventaja narrativa fuerte: puede ser más
interpretable que una MLP común. Podemos presentar ambas rutas sin forzar una falsa pelea:
PPO+MLP como base robusta, Real-KAN como alternativa más novedosa y con mayor resiliencia.

\section*{¿Y la historia? (memoria, atención)}

Probamos LSTM, MLP con historia apilada, y DMLPA/CAN (atención sobre historia). \textbf{Ninguna
mostró una ventaja real sobre PPO+MLP simple.} Esto queda como pregunta abierta para David: puede
que otra forma de codificar la historia sí funcione, pero con lo que probamos esta sesión no
encontramos señal positiva.

\section*{¿Preventivo o reactivo?}

Esta es tu pregunta más interesante. Ya dejamos implementada una auditoría concreta en el repo.
La auditoría no reentrena nada: toma PPO+MLP y Real-KAN ya entrenados, repite los mismos
episodios/semillas, y mide si las acciones aprendidas agregan valor antes o después del evento de
riesgo.

\begin{itemize}
  \item Alinear las acciones de la política con el momento exacto en que empieza cada riesgo, y
        ver si sube turno/despacho \emph{antes} (preventivo) o \emph{después} (reactivo) del
        evento. Esto produce un estudio visual $t=-4,\ldots,+8$ semanas.
  \item Repetir el mismo episodio reemplazando las acciones de una ventana por una acción neutral.
        Si $R_{\mathrm{full}}-R_{\mathrm{reset}}$ es positivo antes del riesgo, hay evidencia de
        prevención; si es positivo después, hay evidencia de reacción/recuperación.
  \item Apagar o mezclar la señal de forecast sólo al evaluar (sin reentrenar): si no cambia el
        desempeño, el modelo no depende causalmente de esa señal para decidir.
  \item Medir sensibilidad por grupos de variables: forecast, fase del régimen, presión aguas
        abajo, backlog y servicio. Para Real-KAN, esto puede complementarse con splines de
        \texttt{pykan}; para CAN/DMLPA, la atención se trata como evidencia auxiliar, no como
        explicación suficiente.
\end{itemize}

La primera corrida pareada de 5 semillas ya terminó y muestra señales descriptivas interesantes:
el valor parece aparecer principalmente después del evento, no antes. Además, al reentrenar sin los
dos forecasts explícitos, la ReT Excel se mantiene prácticamente igual. \textbf{Pero todavía no lo
presentaría como una clasificación final.}
Hicimos un control con tres políticas de referencia (una estática, una diseñada para ser
preventiva y una diseñada para ser reactiva) y la regla automática todavía no separa bien esos
casos. Por tanto, hoy el reclamo defendible es más prudente:
\textbf{hay aprendizaje adaptativo claro, pero la auditoría preventivo/reactivo sigue en
calibración}.

Esto no es una mala noticia; al contrario, nos dice exactamente qué falta para responderte bien.
El siguiente paso es ajustar la auditoría para distinguir dos cosas: (i) si la política sube
turno/despacho antes del riesgo, y (ii) si esa acción preventiva realmente mejora la resiliencia
ReT. El resultado ideal para Garrido sería una política mixta: preventiva ante riesgos comunes y
reactiva ante lo inesperado. Si una siguiente ronda con KAN/CAN logra eso, sería un argumento
fuerte para preferir esa arquitectura.

\section*{Conclusión}

Las dos arquitecturas funcionan y ganan resiliencia real en el contrato 8D completo. PPO+MLP es lo más sólido y eficiente para
el resultado principal del paper. Real-KAN es la alternativa más atractiva si priorizas
resiliencia máxima e interpretabilidad. Queda a tu criterio cuál enfatizar. La auditoría de
prevención queda abierta como la próxima prueba: queremos demostrar no sólo que aprende, sino
\emph{qué} aprende y si ese aprendizaje es preventivo, reactivo o mixto.

\end{document}
